\documentclass[11pt]{article}
\usepackage[margin=1in]{geometry}
\usepackage{amsmath, amssymb}

\title{Exercise 3.5: Modifying the Dynamics for Episodic Tasks}
\author{Sutton \& Barto, Section 3.3}
\date{}

\begin{document}
\maketitle

\section*{Problem}

The equations in Section 3.1 are for the continuing case and need to be
modified (very slightly) to apply to episodic tasks. Show that you know
the modifications needed by giving the modified version of (3.3).

\section*{Background}

Equation (3.3) states that for all $s \in \mathcal{S}$ and $a \in \mathcal{A}(s)$:
\begin{equation}
    \sum_{s' \in \mathcal{S}} \sum_{r \in \mathcal{R}} p(s', r \mid s, a) = 1
    \tag{3.3}
\end{equation}
This says that the dynamics function $p$ defines a valid probability
distribution over next states and rewards, for each state--action pair.

In the continuing case, $\mathcal{S}$ is the set of all states and
$s'$ can be any state in $\mathcal{S}$.

\section*{Solution}

In episodic tasks, we distinguish between:
\begin{itemize}
    \item $\mathcal{S}$ --- the set of all \emph{nonterminal} states
    \item $\mathcal{S}^+$ --- the set of all states including the
          \emph{terminal state}: $\mathcal{S}^+ = \mathcal{S} \cup \{s_{\text{terminal}}\}$
\end{itemize}

From any nonterminal state, the agent can transition to either another
nonterminal state or to the terminal state. Therefore, the sum over
next states must range over $\mathcal{S}^+$ instead of $\mathcal{S}$:

\[
    \boxed{\sum_{s' \in \mathcal{S}^+} \sum_{r \in \mathcal{R}}
    p(s', r \mid s, a) = 1, \quad
    \text{for all } s \in \mathcal{S}, \; a \in \mathcal{A}(s)}
\]

\subsection*{What changes and what doesn't}

\begin{itemize}
    \item The sum over $s'$ now includes $\mathcal{S}^+$ (not $\mathcal{S}$),
          because transitions to the terminal state are possible.
    \item The condition $s \in \mathcal{S}$ (nonterminal) stays the same ---
          there are no actions or transitions \emph{from} the terminal state,
          so the dynamics need only be defined for nonterminal states.
    \item The equation still sums to 1: probability is conserved. Some of
          that probability mass may now go to the terminal state.
\end{itemize}

\end{document}
